\section{Experiments and Results}

This section presents the experimental setup, baseline comparisons, and performance evaluation of the proposed multilingual multimodal cybercrime detection framework.

\subsection{Experimental Setup}

The proposed model is evaluated on the collected dataset of 25000 annotated samples across five languages, namely English, Hindi, Telugu, Tamil, and Malayalam. The dataset is split into training, validation, and test sets in a 70:15:15 ratio. 

Performance is evaluated using standard classification metrics, including accuracy, precision, and F1 score. These metrics are chosen to account for class imbalance and provide a comprehensive evaluation of model performance across multiple cybercrime categories.

\subsection{Baseline Models}

To validate the effectiveness of the proposed approach, comparisons are performed with strong baseline models. These include monolingual and multilingual text-based models, as well as multimodal approaches. Specifically, BERT is used as a monolingual baseline, while mBERT and XLM-R serve as multilingual baselines. Additionally, a vision-language model based on CLIP is used to evaluate multimodal performance.

\subsection{Overall Performance}

Table~\ref{tab:performance} presents the performance comparison. The proposed framework achieves accuracy of 0.892 and F1 of 0.883, outperforming the strongest baseline (XLM-R) by 5.3 F1 points, highlighting the effectiveness of cross-modal attention fusion.

\subsection{Ablation Study}

Table~\ref{tab:ablation} shows component contributions. Removing image modality, multilingual training, or fine-tuning causes substantial performance drops, confirming the necessity of each design choice.

\subsection{Per-Language Analysis}

Table~\ref{tab:language} shows performance across languages. English and Hindi achieve highest scores, while Telugu, Tamil, and Malayalam show slightly lower performance due to data distribution differences.

\subsection{Discussion}

The results confirm that cross-modal fusion effectively captures complementary textual and visual cues. Error analysis reveals lower-performing categories (Dark Web Crimes, Cyber Terrorism) suffer from limited visual evidence and context-dependent patterns. Overall, the framework demonstrates a scalable solution for multilingual cybercrime detection.